\section{Lowest Common Ancestor in a Binary Tree}
\label{sec:trees-lca}

\problemstatement{Given the root of a binary tree and two nodes $p$ and
$q$ (both guaranteed present), return their lowest common ancestor: the
deepest node that has both $p$ and $q$ as descendants (a node is a
descendant of itself).}

\begin{patternbox}
\emph{``Lowest common ancestor'' in a \textbf{general} binary tree} (no BST
ordering to exploit) is a single bottom-up recursion that reports back up
which of $p$ and $q$ each subtree contains. The LCA is the unique node
where the two targets are found on \emph{different sides} -- or which is
itself one target with the other below it.
\end{patternbox}

\subsection*{Understanding the problem carefully}

Walk the tree once. Each recursive call answers: ``does my subtree contain
$p$ or $q$ (or one of them being me)?'' by returning a non-null pointer if
so. If a node finds one target in its left subtree and the other in its
right subtree, that node is the meeting point -- the LCA. If both are found
on the same side, the LCA lies deeper on that side and is propagated up
unchanged.

\begin{ideabox}[The Split Point Is the Answer]
\textsc{LCA}(node): if node is null or equals $p$ or $q$, return node.
Recurse left and right. If \emph{both} calls return non-null, the current
node is where the two targets split -- return it. Otherwise return whichever
side was non-null (or null). The topmost node receiving non-null from both
sides is the LCA.
\end{ideabox}

\subsection*{Algorithm}

\begin{enumerate}
  \item \textsc{LCA}(\textit{node}):
  \begin{enumerate}
    \item If \textit{node} is null or equals $p$ or $q$, return
          \textit{node}.
    \item $L=$\textsc{LCA}(left); $R=$\textsc{LCA}(right).
    \item If $L$ and $R$ are both non-null, return \textit{node}.
    \item Return $L$ if non-null, else $R$.
  \end{enumerate}
\end{enumerate}

\subsection*{Dry run}

\dryrun{Tree: root $3$; left $5$ (children $6,2$); right $1$. Find LCA of
$6$ and $2$.}

\begin{itemize}
  \item At $5$: left call at $6$ returns $6$; right call at $2$ returns
        $2$. Both non-null $\Rightarrow$ $5$ is the LCA; return $5$.
  \item At $3$: left returns $5$, right (subtree $1$) returns null. Only
        left non-null $\Rightarrow$ return $5$.
  \item LCA $=5$.
\end{itemize}

\subsection*{C++ implementation}

\begin{lstlisting}
#include <bits/stdc++.h>
using namespace std;

struct TreeNode {
    int val;
    TreeNode* left;
    TreeNode* right;
    TreeNode(int v) : val(v), left(nullptr), right(nullptr) {}
};

TreeNode* lca(TreeNode* node, TreeNode* p, TreeNode* q) {
    if (node == nullptr || node == p || node == q) return node;
    TreeNode* L = lca(node->left,  p, q);
    TreeNode* R = lca(node->right, p, q);
    if (L != nullptr && R != nullptr) return node;  // split point
    return L != nullptr ? L : R;
}

int main() {
    TreeNode* root = new TreeNode(3);
    root->left = new TreeNode(5);
    root->right = new TreeNode(1);
    TreeNode* p = root->left->left  = new TreeNode(6);
    TreeNode* q = root->left->right = new TreeNode(2);

    cout << "LCA: " << lca(root, p, q)->val << "\n";
    return 0;
}
\end{lstlisting}

\begin{complexitybox}
\textbf{Time Complexity.} $O(n)$ -- a single traversal. \textbf{Space
Complexity.} $O(h)$ for the recursion stack.
\end{complexitybox}

\begin{takeawaybox}
Returning a ``found'' pointer up the recursion and declaring the LCA at
the node where left and right both report success is a one-pass gem. Unlike
the BST version, no value ordering is used -- only the tree's shape --
which is why it works on any binary tree.
\end{takeawaybox}
